% !TeX root = ../main.tex
\chapter{ỨNG DỤNG MINH HOẠ}
\label{app:C}

Để kiểm chứng khả năng vận hành đầu cuối của quy trình, luận văn xây dựng hai công
cụ minh hoạ: một \en{notebook} suy luận trên ảnh đơn lẻ và một ứng dụng web tương
tác.

\section{Suy luận trên ảnh đơn lẻ}

Quy trình suy luận đầy đủ gồm bốn bước nối tiếp trên một ảnh đầu vào:

\begin{enumerate}
    \item \textbf{Phân loại} --- nạp mô hình \EffNet{} đã huấn luyện, trả về lớp
          bệnh dự đoán kèm phân bố xác suất trên năm lớp.
    \item \textbf{Giải thích} --- sinh bản đồ kích hoạt \GradCAM{} cho lớp được dự
          đoán, chồng lên ảnh gốc.
    \item \textbf{Phân đoạn} --- nạp mô hình \UNetpp{} phiên bản V3, trả về bản đồ
          xác suất và mặt nạ nhị phân sau khi áp ngưỡng 0,5.
    \item \textbf{Tổng hợp} --- ghép kết quả thành một bảng năm khung và tính tỉ lệ
          diện tích tổn thương.
\end{enumerate}

\begin{figure}[H]
\centering
\includegraphics[width=\textwidth]{figures/infer_demo.png}
\caption[Kết quả suy luận đầu cuối trên một ảnh]{Kết quả suy luận đầu cuối trên
một ảnh thuộc lớp Algal Leaf Spot. Từ trái sang phải: ảnh đầu vào, phân bố xác
suất năm lớp, bản đồ kích hoạt Grad-CAM chồng lớp, mặt nạ phân đoạn, và ảnh tổng
hợp với đường bao vùng tổn thương.}
\label{fig:infer_demo}
\end{figure}

Với ảnh minh hoạ ở \cref{fig:infer_demo}, mô hình dự đoán lớp Algal Leaf Spot với
độ tin cậy \textbf{99,29\%}. Phân bố xác suất trên năm lớp cho thấy mức độ tách
bạch rất cao: lớp có xác suất lớn thứ hai là Phomopsis Leaf Spot với 0,62\%, các
lớp còn lại đều dưới 0,05\%. Điều này phù hợp với nhận định ở
\cref{sec:ketqua_cls} rằng Algal và Phomopsis là cặp lớp gần nhau nhất --- ngay cả
khi mô hình rất chắc chắn, lớp cạnh tranh duy nhất vẫn là Phomopsis.

\begin{table}[H]
\centering
\caption{Phân bố xác suất dự đoán trên ảnh minh hoạ}
\label{tab:app_infer}
\begin{tabular}{l r}
\toprule
\textbf{Lớp} & \textbf{Xác suất} \\
\midrule
\textbf{Algal Leaf Spot}   & \textbf{0,9929} \\
Phomopsis Leaf Spot        & 0,0062 \\
Leaf Blight                & 0,0005 \\
Healthy Leaf               & 0,0004 \\
Allocaridara Attack        & 0,0000 \\
\bottomrule
\end{tabular}
\end{table}

\section{Ứng dụng web tương tác}

Ứng dụng được xây dựng bằng Streamlit, cho phép người dùng tải lên ảnh lá và nhận
kết quả phân tích tức thời mà không cần biết lập trình.

\subsection{Chức năng}

\begin{itemize}
    \item Tải ảnh lên từ máy người dùng, hỗ trợ định dạng JPG và PNG.
    \item Chọn phiên bản mô hình phân đoạn trong ba lựa chọn V1, V2, V3 --- phục
          vụ so sánh trực quan giữa các phiên bản.
    \item Điều chỉnh ngưỡng nhị phân hoá bằng thanh trượt trong khoảng 0,1--0,9,
          cho phép quan sát ảnh hưởng của ngưỡng lên diện tích mặt nạ.
    \item Bật hoặc tắt lớp hiển thị bản đồ kích hoạt.
    \item Hiển thị đồng thời năm khung: ảnh gốc, biểu đồ xác suất, bản đồ kích
          hoạt, mặt nạ phân đoạn, và ảnh tổng hợp có đường bao.
    \item Hiển thị bản đồ xác suất thô và thông tin điểm kiểm tra đang dùng.
\end{itemize}

\subsection{Kiến trúc kỹ thuật}

\begin{table}[H]
\centering
\caption{Thành phần kỹ thuật của ứng dụng minh hoạ}
\label{tab:app_streamlit}
\begin{tabular}{L{4.6cm} L{8.4cm}}
\toprule
\textbf{Thành phần} & \textbf{Mô tả} \\
\midrule
Giao diện          & Streamlit, bố cục rộng, thanh điều khiển bên trái \\
Biểu đồ            & Plotly, biểu đồ cột ngang cho phân bố xác suất \\
Nạp mô hình        & Có bộ nhớ đệm, mô hình chỉ nạp một lần cho mỗi phiên làm việc \\
Thiết bị tính toán & Tự phát hiện GPU, chuyển về CPU nếu không có \\
Bảng màu           & Okabe--Ito, an toàn với người mù màu \\
Đầu ra định lượng  & Lớp dự đoán, độ tin cậy, tỉ lệ diện tích tổn thương \\
\bottomrule
\end{tabular}
\end{table}

Việc cho phép người dùng chuyển đổi giữa ba phiên bản mô hình ngay trên giao diện
có giá trị sư phạm: nó cho thấy trực quan điều mà \cref{sec:caveat} phân tích bằng
số liệu --- ba phiên bản cho ra ba kiểu mặt nạ khác nhau về hình dạng và diện tích,
và không thể kết luận phiên bản nào ``đúng'' hơn nếu không có nhãn chuẩn.

\subsection{Hạn chế của ứng dụng}

Ứng dụng chỉ nhằm mục đích minh hoạ, chưa phải sản phẩm triển khai thực tế. Bốn
hạn chế chính:

\begin{itemize}
    \item Chưa hiệu chỉnh sai lệch hệ thống của mô hình --- như đã phân tích ở
          \cref{sec:ablation}, mô hình phiên bản V3 có xu hướng ước lượng thừa
          diện tích tổn thương.
    \item Chỉ xử lý ảnh chụp lá đơn lẻ trên nền tương đối sạch, chưa xử lý được
          ảnh chụp toàn tán.
    \item Tỉ lệ diện tích tổn thương hiển thị là tỉ lệ trên \emph{toàn khung ảnh},
          không phải trên diện tích lá --- hai đại lượng chỉ trùng nhau khi lá
          chiếm trọn khung.
    \item Chưa có cơ chế cảnh báo khi ảnh đầu vào nằm ngoài phân bố dữ liệu huấn
          luyện, chẳng hạn ảnh không phải lá sầu riêng.
\end{itemize}
